% Self-contained mathematical fragment; requires amsmath, amssymb, amsthm.
% Suggested bibliography keys: Lagarias2007, BMOR2021, NakamuraSuzuki2023.

\section{A universal process built from zero ordinates}

This section separates two facts. The probabilistic theorem below follows
from the classical zero-counting law alone. RH or Dirichlet GRH identifies
the resulting circle generator with the usual Li coefficients. The repository supplies the established zero-line theorem needed for
that identification.

Let $\mathcal Z$ be a locally finite multiset of real numbers, with
multiplicities, and put
\[
 N(T)=\#\{\gamma\in\mathcal Z:|\gamma|\leq T\}.
\]
Assume that, for constants $d>0$ and $B\in\mathbb R$,
\begin{equation}\label{eq:generic-zero-count}
 N(T)=\frac{d}{\pi}T\log T+\frac{B}{\pi}T+O(\log(T+2))
 \qquad(T\longrightarrow\infty).
\end{equation}
Define
\[
 \theta(\gamma)=
 \begin{cases}2\arctan(1/(2\gamma)),&\gamma\ne0,\\
 \pi,&\gamma=0,
 \end{cases}
 \qquad
 \nu=\frac12\sum_{\gamma\in\mathcal Z}
       (\delta_{\theta(\gamma)}+\delta_{-\theta(\gamma)}).
\]
In particular, $m$ occurrences of $\gamma=0$ contribute mass $m/2$ at
each of $\pi$ and $-\pi$. Write
\begin{equation}\label{eq:jump-symbol}
 \Psi(u)=\int_{\mathbb R}(1-\cos(uy))\,\nu(dy)
       =\sum_{\gamma\in\mathcal Z}(1-\cos(u\theta(\gamma))).
\end{equation}

\begin{lemma}[Existence and elementary properties]\label{lem:jump-exists}
The measure $\nu$ is a symmetric L\'evy measure supported in $[-\pi,\pi]$.
It satisfies
\[
 \nu(\mathbb R)=\infty,\qquad
 \int |y|\,\nu(dy)=\infty,\qquad
 0<V:=\int y^2\,\nu(dy)<\infty.
\]
There is a symmetric L\'evy process $(X_t)_{t\geq0}$ with
\[
 \mathbb E e^{iuX_t}=e^{-t\Psi(u)}.
\]
Its Gaussian component is zero, it has infinitely many jumps and infinite
variation, and $\mathbb E X_t^2=tV$.
\end{lemma}

\begin{proof}
The counting law and partial summation show that
$\sum_{|\gamma|>1}|\gamma|^{-2}<\infty$ whereas
$\sum_{|\gamma|>1}|\gamma|^{-1}=\infty$.
Since $|\theta(\gamma)|\sim |\gamma|^{-1}$, all assertions about the
measure follow. L\'evy--Khintchine gives the process, with zero Gaussian
coefficient and zero symmetric drift. The standard variation criterion
for a pure-jump process gives infinite variation. Differentiating the
characteristic function twice at zero gives the variance identity.
\end{proof}

The construction can also be understood by retaining finitely many
ordinates, allowing independent Poisson jumps of each indicated size,
and taking the symmetric limit. The finite second-moment tail makes this
limit well defined; infinitely many small jumps remain in the limit.

\begin{lemma}[The symbol at large frequency]\label{lem:continuous-symbol}
Put $A=d/2$ and $\kappa=\gamma_{\!E}+B/d$, where $\gamma_{\!E}$ is
Euler's constant. Then
\begin{equation}\label{eq:symbol-asymptotic}
 \Psi(u)=A|u|(\log|u|+\kappa)
          +O(\sqrt{|u|}\log(|u|+2))
 \qquad(|u|\longrightarrow\infty).
\end{equation}
\end{lemma}

\begin{proof}
By evenness take $u\geq2$ and set $T=\sqrt u$. The contribution of
$|\gamma|\leq T$ is $O(T\log(T+2))$. On the remaining ordinates,
\[
 \theta(\gamma)=\gamma^{-1}+O(|\gamma|^{-3}),
\]
so the Lipschitz bound for cosine makes the replacement by
$1-\cos(u/\gamma)$ cost at most
\[
 O\left(u\int_T^\infty v^{-3}\,dN(v)\right)=O(\log u).
\]
Let $M(v)=(d v\log v+Bv)/\pi$ and $R(v)=N(v)-M(v)$.
For $g(v)=1-\cos(u/v)$, integration by parts gives
\[
 \int_T^\infty g(v)\,dR(v)
 =-g(T)R(T)-\int_T^\infty R(v)g'(v)\,dv.
\]
The boundary term at infinity vanishes since $g(v)=O(u^2/v^2)$.
As $|g'(v)|\leq u/v^2$, the displayed expression is
\[
 O\left(\log(T+2)+u\int_T^\infty
              \frac{\log(v+2)}{v^2}\,dv\right)
 =O(\sqrt u\log u).
\]
Extending the main integral from $T$ down to zero costs
$O(T\log(T+2))$. Consequently,
\begin{align*}
 \Psi(u)
 &=\frac1\pi\int_0^\infty
     (1-\cos(u/v))(d\log v+d+B)\,dv+O(\sqrt u\log u)\\
 &=\frac u\pi\left[(d\log u+d+B)I_0-dI_1\right]
        +O(\sqrt u\log u),
\end{align*}
where
\[
 I_0=\int_0^\infty\frac{1-\cos y}{y^2}\,dy=\frac\pi2,
 \qquad
 I_1=\int_0^\infty\frac{(1-\cos y)\log y}{y^2}\,dy
     =\frac\pi2(1-\gamma_{\!E}).
\]
For completeness these evaluations follow by differentiating, at
$\alpha=1$, the absolutely convergent identity
\[
 \int_0^\infty(1-\cos y)y^{-1-\alpha}\,dy
 =\frac{\pi}{2\Gamma(1+\alpha)\sin(\pi\alpha/2)},
 \qquad 0<\alpha<2.
\]
One proof of the identity inserts the Laplace representation of
$y^{-1-\alpha}$ and uses Tonelli's theorem and
$\int_0^\infty e^{-ty}(1-\cos y)\,dy=1/[t(1+t^2)]$.
Substitution gives \eqref{eq:symbol-asymptotic}.
\end{proof}

\section{The all-order density expansion}

Let $W_0$ denote the positive real branch of Lambert's function,
characterized by $W_0(z)e^{W_0(z)}=z$ for $z>0$. Define
\begin{equation}\label{eq:generic-scale}
 w(t)=W_0\left(\frac{e^\kappa}{At}\right),\qquad
 a(t)=Atw(t).
\end{equation}
The exact identities
\begin{equation}\label{eq:scale-identities}
 a=e^{\kappa-w},\qquad \log(1/a)+\kappa=w,
 \qquad t=\frac{a}{Aw}
\end{equation}
are the reason to use Lambert $W$ instead of truncating the scale itself.

For $k\geq0$ set
\begin{equation}\label{eq:universal-profiles}
 Q_k(x)=\frac{(-1)^k}{\pi k!}\operatorname{Re}
 \left.
 \frac{\partial^k}{\partial z^k}
       \bigl[\Gamma(z)(1-ix)^{-z}\bigr]
 \right|_{z=k+1}.
\end{equation}
The logarithm implicit in the complex power is the principal logarithm;
its argument always has positive real part.

\begin{theorem}[Universal expansion at every logarithmic order]
\label{thm:all-orders}
For every $t>0$, $X_t$ has a density $f_t$ which extends to an entire
function of its spatial variable. For every pair of fixed integers
$K,r\geq0$,
\begin{equation}\label{eq:all-order-bound}
 \sup_{x\in\mathbb R}
 \left|
 \frac{d^r}{dx^r}
 \left[a(t)f_t(a(t)x)-\sum_{k=0}^K\frac{Q_k(x)}{w(t)^k}\right]
 \right|
 =O_{K,r}\bigl(w(t)^{-K-1}+\sqrt{a(t)}\bigr)
\end{equation}
as $t\downarrow0$. In particular the error is $O_{K,r}(w^{-K-1})$.
The constants may depend on the multiset and its counting estimate.
\end{theorem}

\begin{proof}
The large-frequency estimate makes
$\exp(-t\Psi(u)+c|u|)$ integrable for every $t>0$ and $c>0$.
Fourier inversion therefore gives an entire density.
Put $a=a(t)$, $w=w(t)$, and for $s>0$ define
\[
 F_w(s)=s+\frac{s\log s}{w}.
\]
The scale identities yield
\[
 tA\frac{s}{a}\left(\log\frac{s}{a}+\kappa\right)=F_w(s).
\]
We first establish, for each fixed $r\geq0$,
\begin{equation}\label{eq:fourier-comparison}
 \int_0^\infty s^r
  \left|e^{-t\Psi(s/a)}-e^{-F_w(s)}\right|\,ds
 =O_r(\sqrt a).
\end{equation}
On $0<s<\sqrt a$, the first exponential is at most one. The negative
part of $F_w$ is at most $e^{-w-1}/w$, so the second exponential is
uniformly bounded as well. This interval contributes $O_r(\sqrt a)$.
On $s\geq\sqrt a$, both exponents are at least $cs$ for some fixed
$c>0$, once $t$ is small enough. Indeed
$F_w(s)/s\geq 1/2+\kappa/(2w)$, and the remainder in
\eqref{eq:symbol-asymptotic} is smaller than a fixed fraction of its
main term on $s/a\geq a^{-1/2}$. The exponential mean-value bound now gives
\[
 \left|e^{-t\Psi(s/a)}-e^{-F_w(s)}\right|
 \leq C t\sqrt{s/a}\log(2+s/a)e^{-cs}.
\]
Since $t=a/(Aw)$ and
$\log(2+s/a)\leq w+C_\kappa+\log(2+s)$, its integral against $s^r$
is $O_r(\sqrt a)$, proving \eqref{eq:fourier-comparison}.

Expand $e^{-s\log s/w}$ through degree $K$. For $s\geq1$, Taylor's
integral remainder for $e^{-z}$, $z\geq0$, is bounded by
$z^{K+1}/(K+1)!$. For $0<s<1$, use $|s\log s|\leq1/e$ and the
same estimate with the harmless factor $e^{1/(ew)}$. Multiplication by
$e^{-s}s^r$ and integration show that
\begin{equation}\label{eq:taylor-comparison}
 \int_0^\infty s^r e^{-s}
 \left|e^{-s\log s/w}
  -\sum_{k=0}^K\frac{(-s\log s)^k}{k!w^k}\right|\,ds
 =O_{K,r}(w^{-K-1}).
\end{equation}
The sum of \eqref{eq:fourier-comparison} and
\eqref{eq:taylor-comparison} is an $L^1$ bound for the Fourier
transforms, with every fixed polynomial weight. Fourier inversion and
differentiation prove the uniform claim. Finally,
\[
 \int_0^\infty s^k(\log s)^k e^{-(1-ix)s}\,ds
 =\left.\partial_z^k
       [\Gamma(z)(1-ix)^{-z}]\right|_{z=k+1},
\]
which identifies the profiles. As $a=e^{\kappa-w}$, $\sqrt a$ is
smaller than every fixed inverse power of $w$.
\end{proof}

The first three terms have the explicit form
\begin{align}
 Q_0(x)&=\frac1{\pi(1+x^2)},\label{eq:profile-zero}\\
 Q_1(x)&=-\frac1\pi\operatorname{Re}
       \frac{1-\gamma_{\!E}-\log(1-ix)}{(1-ix)^2},
       \label{eq:profile-one}\\
 Q_2(x)&=\frac1\pi\operatorname{Re}
       \frac{(\frac32-\gamma_{\!E}-\log(1-ix))^2+\frac{\pi^2}6-\frac54}
            {(1-ix)^3}.
       \label{eq:profile-two}
\end{align}
Thus the leading shape is the standard Cauchy density, and the theorem
provides every successive correction rather than only convergence in law.

\begin{corollary}[Universality beyond every logarithmic order]
If two multisets satisfy \eqref{eq:generic-zero-count} with the same
$d,B$, their scaled densities differ by $O_r(\sqrt a)$ in
$C_b^r(\mathbb R)$. In particular they have the same expansion to every
fixed order in $1/\log(1/t)$.
\end{corollary}

\begin{proof}
Both scaled characteristic functions differ from the same model
$e^{-F_w(|s|)}$ by $O_r(\sqrt a)$ in the weighted $L^1$ norm
\eqref{eq:fourier-comparison}. Apply Fourier inversion directly; this
is stronger than subtracting two fixed-order asymptotic expansions.
\end{proof}

\section{Dirichlet characters and Li coefficients}

For zeta, use the normalization
\[\xi(s)=\tfrac12s(s-1)\pi^{-s/2}\Gamma(s/2)\zeta(s),\qquad
\lambda_n=\frac1{(n-1)!}\left.\frac{d^n}{ds^n}
[s^{n-1}\log\xi(s)]\right|_{s=1}.\]
The equivalent zero-sum formula and its Dirichlet generalization
are used below~\cite{Li1997,BombieriLagarias,Lagarias2007}.

Take $\mathcal Z$ to be the ordinates of all nontrivial zeros of a fixed
primitive Dirichlet $L$-function of conductor $q>1$, counted with
multiplicity and both signs. The classical counting law, in the explicit
form of Bennett--Martin--O'Bryant--Rechnitzer~\cite{BMOR2021}, gives
\[
 N_\chi(T)=\frac T\pi\log\frac{qT}{2\pi e}
              +O_\chi(\log(T+2)).
\]
The zeta function is included by taking $q=1$ and its usual
Riemann--von Mangoldt formula. Thus $d=1$,
$B=\log(q/(2\pi))-1$, and the scale is
\begin{equation}\label{eq:dirichlet-scale}
 w_q(t)=W_0\left(\frac{q e^{\gamma_{\!E}-1}}{\pi t}\right),
 \qquad a_q(t)=\frac t2 w_q(t).
\end{equation}
All preceding density statements are unconditional for these ordinate
processes. All primitive characters of a fixed conductor have the same
expansion to every logarithmic order. Constants are for fixed characters;
no uniformity in $q\to\infty$ is asserted.

Under the relevant RH, each zero $\rho=1/2+i\gamma$ satisfies
$1-1/\rho=e^{i\theta(\gamma)}$ (with $e^{i\pi}=-1$ at a central
zero). Hence the usual zero-sum Li coefficient obeys
\begin{equation}\label{eq:li-identification}
 \operatorname{Re}\lambda_\chi(n)
 =\sum_\rho\operatorname{Re}
       \left[1-\left(1-\frac1\rho\right)^n\right]
 =\Psi_\chi(n).
\end{equation}
The real-part sum is absolutely convergent. For real characters the Li
coefficients themselves are real; for nonreal characters the real part
in \eqref{eq:li-identification} is essential.

The link with Li coefficients and positivity has substantial precedent~\cite{Lagarias2007,NakamuraSuzuki2023}.
The all-order normalized density expansion, rather than the existence
of positive kernels alone, is the proposed contribution here.

\section{The circle kernel and its trace}

Reduce $X_t$ modulo $2\pi$ and let $p_t$ be its density with respect to
ordinary Lebesgue measure on $[-\pi,\pi]$. Its Fourier series is
\begin{equation}\label{eq:circle-kernel}
 p_t(y)=\frac1{2\pi}\sum_{n\in\mathbb Z}e^{-t\Psi(n)}e^{iny}.
\end{equation}
The series converges locally uniformly for complex $y$, so the density
is entire and periodic. Under RH or GRH the associated nonnegative generator has eigenvalues
$\operatorname{Re}\lambda_\chi(|n|)$, with $\lambda_\chi(0)=0$.

\begin{proposition}[Local circle expansion and heat trace]
For every fixed $R>0$ and integers $K,r\geq0$,
\[
 \sup_{|x|\leq R}\left|
 \frac{d^r}{dx^r}
 \left[ap_t(ax)-\sum_{k=0}^K\frac{Q_k(x)}{w^k}\right]
 \right|=O_{R,K,r}(w^{-K-1}+\sqrt a).
\]
If $\mathcal A$ is the circle generator, then
\begin{align*}
 \operatorname{Tr}(e^{-t\mathcal A})
 &=2\pi p_t(0)\\
 &=\frac2a\left[
 1-\frac{1-\gamma_{\!E}}w
 +\frac{(\frac32-\gamma_{\!E})^2+\frac{\pi^2}6-\frac54}{w^2}
 +O(w^{-3})\right].
\end{align*}
For a primitive Dirichlet character the leading trace is
$4/(t w_q(t))$.
\end{proposition}

\begin{proof}
The scaled sum \eqref{eq:circle-kernel} has mesh $a$ in frequency.
Splitting at $an=\sqrt a$ and applying the same mean-value estimates
as in \eqref{eq:fourier-comparison} bounds its difference from the
model Riemann sum by $O_r(\sqrt a)$. The model function
$s^r e^{-F_w(s)}e^{ixs}$ has uniformly integrable first derivative for
$|x|\leq R$; the only possible singularity is an integrable logarithm
at zero. The Riemann-sum error is consequently $O_{R,r}(a)$.
The Taylor calculation is unchanged. Setting $x=0$ gives the trace.
\end{proof}

Uniformity over all real $x$ in the circle assertion would be false:
periodicity produces returning peaks at distances $2\pi/a$. The
whole-line theorem has no such restriction.

\section{A moment transition and two different time scales}

The Cauchy limit does not contradict the finite variance of $X_t$.
The family loses uniform integrability after the short-time
normalization. The following formulas quantify the change.

\begin{theorem}[Moment transition]\label{thm:moment-transition}
Let $L=\log(1/t)$ and retain $A=d/2$. As $t\downarrow0$,
\begin{align*}
 \mathbb E|X_t|^p
 &\sim \frac{a(t)^p}{\cos(\pi p/2)},&&0<p<1,\\
 \mathbb E|X_t|
 &\sim \frac A\pi tL^2
       =\frac d{2\pi}tL^2,&&p=1,\\
 \mathbb E|X_t|^p
 &\sim t\int |y|^p\nu(dy),&&1<p<2,\\
 \mathbb E X_t^2&=tV.&&
\end{align*}
Moreover
\[
 \frac{X_t}{a(t)}\ \Longrightarrow\ \mathrm{Cauchy}(0,1)
 \quad(t\downarrow0),\qquad
 \frac{X_t}{\sqrt{tV}}\ \Longrightarrow\ N(0,1)
 \quad(t\longrightarrow\infty).
\]
\end{theorem}

\begin{proof}
For $0<p<2$ the fractional-moment identity is
\[
 \mathbb E|X_t|^p
 =c_p\int_0^\infty(1-e^{-t\Psi(u)})u^{-p-1}\,du,
 \qquad c_p=\frac{2\Gamma(1+p)\sin(\pi p/2)}\pi.
\]
It follows by Tonelli from the cosine integral evaluated in
Lemma~\ref{lem:continuous-symbol}. For $0<p<1$ put $u=s/a$.
For a sufficiently large fixed $U_0$, the interval $s\leq aU_0$
contributes $O(t a^{-p})=o(1)$ to the normalized moment, using
$\Psi(u)\leq Vu^2/2$. On $aU_0\leq s\leq1$, the symbol bound gives
$t\Psi(s/a)\leq Cs$; on $s\geq1$ use $1-e^{-t\Psi}\leq1$.
Dominated convergence then gives
\[
 \lim_{t\downarrow0}\frac{\mathbb E|X_t|^p}{a^p}
 =c_p\int_0^\infty(1-e^{-s})s^{-p-1}\,ds
 =\frac1{\cos(\pi p/2)}.
\]
For $1<p<2$, divide the unscaled identity by $t$ and use
$(1-e^{-t\Psi})/t\leq\Psi$. The bound
$\Psi(u)=O(u^2)$ at zero and \eqref{eq:symbol-asymptotic} at infinity
provide an integrable majorant. Tonelli identifies the limit with
$\int|y|^p\nu(dy)$.

At $p=1$, the part $0<u<1$ is $O(t)$. For an upper bound split at
$U=1/t$, use $1-e^{-z}\leq\min(z,1)$, and obtain
\[
 \mathbb E|X_t|
 \leq \frac{2t}\pi\int_1^{1/t}\frac{\Psi(u)}{u^2}\,du+O(t)
 =\frac A\pi tL^2+O(tL).
\]
For a lower bound retain only $1\leq u\leq1/(tL^2)$ and use
$1-e^{-z}\geq z-z^2/2$. The main term is
$\frac A\pi t(L-2\log L)^2+O(tL)$; the quadratic error is $O(t)$,
since $t^2\int_1^{1/(tL^2)}\Psi(u)^2u^{-2}\,du=O(t)$.
These two bounds give the claimed equivalent. The variance identity
was proved earlier.

The short-time weak limit follows from
$t\Psi(s/a)\to|s|$. For long times, dominated convergence in the
L\'evy integral gives $\Psi(u)=Vu^2/2+o(u^2)$ as $u\to0$.
Consequently $t\Psi(s/\sqrt{tV})\to s^2/2$, which proves the
Gaussian limit by the continuity theorem.
\end{proof}

